% GENERATED FILE. Edit canonical lessons, then run npm run generate:books.
% canonical-source-hash: fnv1a64:266eb130d7544e7a
% canonical-lessons: RU-C353-ushcherb, RU-C353-prichinyat, RU-C353-vredit, RU-C353-stradat, RU-C353-ukhudshat

\chapter{Damage and harm}
\label{ch:ru-damage-and-harm}
\hlchaptermodality{353}

\begin{chapteropening}
\textbf{By the end of this chapter:} \emph{I can speak about damage, say what causes harm, what does damage and what makes things worse, and who suffers.}

Complete the last lesson of chapter 353: I can speak about damage, say what causes harm, what does damage and what makes things worse, and who suffers.
\end{chapteropening}

\section[\texorpdfstring{\ru{ущерб} (ushchérb)}{ущерб (ushchérb)}]{\ru{ущерб} (ushchérb) — damage, loss}

\label{lesson:RU-C353-ushcherb}



\begin{quote}
Before the new one: say the Russian for an excuse, then the Russian for responsible.
\end{quote}

\subsection*{You'll want to know: \ru{ущерб}}
\textbf{\ru{ущерб}} (\emph{ushchérb}) — \textquotedblleft{}damage, loss\textquotedblright{}.

\begin{quote}
\ru{Пожар} \ru{причинил} \ru{большой} \ru{ущерб}.
\end{quote}

\begin{quote}
\emph{Pozhár prichiníl bolshóy ushchérb.}
\end{quote}

\begin{quote}
The fire caused great damage.
\end{quote}

Masculine. \emph{\ru{Нанести} \ru{ущерб}} — to cause damage. \emph{\ru{Луна} \ru{на} \ru{ущербе}} — the moon is waning.

\begin{cousinweb}
\textbf{\ru{ущерб}} is usually traced to \emph{\ru{у}-} and the root of \emph{\ru{щерба}}, \emph{\ru{щербина}} \textquotedblleft{}a notch, a gap\textquotedblright{}: damage is a piece broken out of something.
\end{cousinweb}

\begin{grammarlens}[title={What you've built}]
One more word for telling, explaining and describing at B1.
\end{grammarlens}

\subsection*{Your turn}
\begin{itemize}
\raggedright
  \item \emph{Say it:} \emph{ushchérb}
  \item \emph{Say it:} \emph{ushchérb}, once more
  \item \emph{Say it:} a sentence of your own with \emph{ushchérb}
  \item \emph{Recall:} say the Russian for an excuse, then the Russian for responsible, then say \emph{ushchérb} again
\end{itemize}

\begin{tcolorbox}[breakable,colback=teal!4,colframe=teal!35!black,title={Before you move on}]
What does \ru{ущерб} mean? (\textquotedblleft{}Damage, loss\textquotedblright{}.) Say it once more.
\end{tcolorbox}
\section[\texorpdfstring{\ru{причинять} (prichinyát)}{причинять (prichinyát)}]{\ru{причинять} (prichinyát) — to cause (harm)}

\label{lesson:RU-C353-prichinyat}



\begin{quote}
Before the new one: say the Russian for responsible, then the Russian for damage.
\end{quote}

\subsection*{You'll want to know: \ru{причинять}}
\textbf{\ru{причинять}} (\emph{prichinyát}) — \textquotedblleft{}to cause (harm)\textquotedblright{}.

\begin{quote}
\ru{Я} \ru{не} \ru{хотел} \ru{причинять} \ru{тебе} \ru{боль}.
\end{quote}

\begin{quote}
\emph{Ya ne khotél prichinyát tebé bol.}
\end{quote}

\begin{quote}
I did not want to cause you pain.
\end{quote}

Imperfective; the perfective is \emph{\ru{причинить}}. It goes with bad things: \emph{\ru{причинять} \ru{боль}, \ru{вред}, \ru{ущерб}}.

\begin{cousinweb}
\textbf{\ru{причинять}} is the verb made from \emph{\ru{причина}} \textquotedblleft{}a cause\textquotedblright{}, which is usually traced to \emph{\ru{при}-} and \emph{\ru{чин}} \textquotedblleft{}order, rank\textquotedblright{}, the root of \emph{\ru{чинить}} \textquotedblleft{}to make, to mend\textquotedblright{}.
\end{cousinweb}

\begin{grammarlens}[title={What you've built}]
One more word for telling, explaining and describing at B1.
\end{grammarlens}

\subsection*{Your turn}
\begin{itemize}
\raggedright
  \item \emph{Say it:} \emph{prichinyát}
  \item \emph{Say it:} \emph{prichinyát}, once more
  \item \emph{Say it:} a sentence of your own with \emph{prichinyát}
  \item \emph{Recall:} say the Russian for responsible, then the Russian for damage, then say \emph{prichinyát} again
\end{itemize}

\begin{tcolorbox}[breakable,colback=teal!4,colframe=teal!35!black,title={Before you move on}]
What does \ru{причинять} mean? (\textquotedblleft{}To cause (harm)\textquotedblright{}.) Say it once more.
\end{tcolorbox}
\section[\texorpdfstring{\ru{вредить} (vredít)}{вредить (vredít)}]{\ru{вредить} (vredít) — to harm, to do damage to}

\label{lesson:RU-C353-vredit}



\begin{quote}
Before the new one: say the Russian for damage, then the Russian for to cause.
\end{quote}

\subsection*{You'll want to know: \ru{вредить}}
\textbf{\ru{вредить}} (\emph{vredít}) — \textquotedblleft{}to harm, to do damage to\textquotedblright{}.

\begin{quote}
\ru{Сладкое} \ru{вредит} \ru{зубам}.
\end{quote}

\begin{quote}
\emph{Sládkoye vredít zubám.}
\end{quote}

\begin{quote}
Sweet things harm your teeth.
\end{quote}

Imperfective; present \emph{\ru{врежу}, \ru{вредишь}}; perfective \emph{\ru{навредить}}. It takes the dative.

\begin{cousinweb}
\textbf{\ru{вредить}} is the verb made from \emph{\ru{вред}} \textquotedblleft{}harm\textquotedblright{}, which Russian took from Church Slavonic \emph{vrědŭ}, from Proto-Slavic \emph{*verdŭ} (EtymDB).
\end{cousinweb}

\begin{grammarlens}[title={What you've built}]
One more word for telling, explaining and describing at B1.
\end{grammarlens}

\subsection*{Your turn}
\begin{itemize}
\raggedright
  \item \emph{Say it:} \emph{vredít}
  \item \emph{Say it:} \emph{vredít}, once more
  \item \emph{Say it:} a sentence of your own with \emph{vredít}
  \item \emph{Recall:} say the Russian for damage, then the Russian for to cause, then say \emph{vredít} again
\end{itemize}

\begin{tcolorbox}[breakable,colback=teal!4,colframe=teal!35!black,title={Before you move on}]
What does \ru{вредить} mean? (\textquotedblleft{}To harm, to do damage to\textquotedblright{}.) Say it once more.
\end{tcolorbox}
\section[\texorpdfstring{\ru{страдать} (stradát)}{страдать (stradát)}]{\ru{страдать} (stradát) — to suffer}

\label{lesson:RU-C353-stradat}



\begin{quote}
Before the new one: say the Russian for to cause, then the Russian for to harm.
\end{quote}

\subsection*{You'll want to know: \ru{страдать}}
\textbf{\ru{страдать}} (\emph{stradát}) — \textquotedblleft{}to suffer\textquotedblright{}.

\begin{quote}
\ru{Он} \ru{страдает} \ru{от} \ru{головной} \ru{боли}.
\end{quote}

\begin{quote}
\emph{On stradáyet ot golovnóy bóli.}
\end{quote}

\begin{quote}
He suffers from headaches.
\end{quote}

Imperfective. \emph{\ru{Страдать} \ru{от}} + genitive (from a cause), \emph{\ru{страдать}} + instrumental (from an illness).

\begin{cousinweb}
\textbf{\ru{страдать}} belongs with \emph{\ru{страда}} \textquotedblleft{}hard farm work, harvest time\textquotedblright{} (EtymDB lists both, and glosses the verb \textquotedblleft{}to suffer, to drudge\textquotedblright{}). Toil and suffering sit side by side in the one word.
\end{cousinweb}

\begin{grammarlens}[title={What you've built}]
One more word for telling, explaining and describing at B1.
\end{grammarlens}

\subsection*{Your turn}
\begin{itemize}
\raggedright
  \item \emph{Say it:} \emph{stradát}
  \item \emph{Say it:} \emph{stradát}, once more
  \item \emph{Say it:} a sentence of your own with \emph{stradát}
  \item \emph{Recall:} say the Russian for to cause, then the Russian for to harm, then say \emph{stradát} again
\end{itemize}

\begin{tcolorbox}[breakable,colback=teal!4,colframe=teal!35!black,title={Before you move on}]
What does \ru{страдать} mean? (\textquotedblleft{}To suffer\textquotedblright{}.) Say it once more.
\end{tcolorbox}
\section[\texorpdfstring{\ru{ухудшать} (ukhudshát)}{ухудшать (ukhudshát)}]{\ru{ухудшать} (ukhudshát) — to make worse}

\label{lesson:RU-C353-ukhudshat}



\begin{quote}
Before the new one: say the Russian for to harm, then the Russian for to suffer.
\end{quote}

\subsection*{You'll want to know: \ru{ухудшать}}
\textbf{\ru{ухудшать}} (\emph{ukhudshát}) — \textquotedblleft{}to make worse\textquotedblright{}.

\begin{quote}
\ru{Стресс} \ru{ухудшает} \ru{сон}.
\end{quote}

\begin{quote}
\emph{Stress ukhudsháyet son.}
\end{quote}

\begin{quote}
Stress makes your sleep worse.
\end{quote}

Imperfective; perfective \emph{\ru{ухудшить}}. \emph{\ru{Ухудшаться}} — to get worse.

\begin{cousinweb}
\textbf{\ru{ухудшать}} is built on \emph{\ru{худший}} \textquotedblleft{}worse\textquotedblright{}, from \emph{\ru{худой}} \textquotedblleft{}thin, poor, bad\textquotedblright{}, inherited from Proto-Slavic \emph{*xudŭ} (EtymDB).
\end{cousinweb}

\begin{grammarlens}[title={What you've built}]
One more word for telling, explaining and describing at B1.
\end{grammarlens}

\subsection*{Your turn}
\begin{itemize}
\raggedright
  \item \emph{Say it:} \emph{ukhudshát}
  \item \emph{Say it:} \emph{ukhudshát}, once more
  \item \emph{Say it:} a sentence of your own with \emph{ukhudshát}
  \item \emph{Recall:} say the Russian for to harm, then the Russian for to suffer, then say \emph{ukhudshát} again
\end{itemize}

\begin{tcolorbox}[breakable,colback=teal!4,colframe=teal!35!black,title={Before you move on}]
What does \ru{ухудшать} mean? (\textquotedblleft{}To make worse\textquotedblright{}.) Say it once more.
\end{tcolorbox}
